\begin{table}[t]
    \centering
    \resizebox{0.85\textwidth}{!}{%
    \begin{tabular}{@{}llccccc@{}}
        \toprule
        Arm & Family & \bag in demented & \bag in controls & ROC-AUC [95\% CI] & Cohen's $d$ & $p$ \\
        \midrule
        \tabpfn & tabular & {\bf $+4.26$\,y} & 0.00 & {\bf 0.701} [0.646, 0.757] & {\bf 0.730} & $<10^{-4}$ \\
        \logreg & tabular & $+3.90$\,y & 0.00 & 0.691 [0.633, 0.749] & 0.660 & $<10^{-4}$ \\
        \xgb    & tabular & $+4.10$\,y & 0.00 & 0.683 [0.628, 0.737] & 0.664 & $<10^{-4}$ \\
        \midrule
        \cotarm    & llm     & $+0.15$\,y & 0.00 & 0.490 [0.435, 0.547] & 0.015 & 0.670 \\
        \direct & llm     & $-0.20$\,y & 0.00 & 0.470 [0.408, 0.535] & $-0.019$ & 0.202 \\
        \bottomrule
    \end{tabular}
    }
    \caption{
        The clinical endpoint. Normative models are fitted on cognitively normal scans only and the
        regression-to-the-mean bias is removed within controls, so the control gap is zero by
        construction. The conventional models reproduce the textbook finding of roughly four years
        of apparent brain ageing in dementia; the \cotarm-derived gap carries no disease signal at all
        (AUC 0.490, $p=0.670$). The chains are articulate about atrophy and ageing, and the quantity
        they produce is diagnostically empty.
    }
    \label{tab:brainage}
\end{table}
